\begin{helper}
Let \(\delta\) be a real number with \(0<\delta<1/10\).  There is a real
threshold \(X_0\) such that, whenever \(x\ge X_0\), there is an integer
\(m\ge1\).  If \(q=2^m\), then the concrete field
\(\mathrm{GaloisField}(2,m)\) has \(q\) elements and
\[
  q\le \frac{\delta}{100}\frac{x}{\log x}\le 2q.
\]
Moreover
\[
  x^{1-\delta/2}\le q,\qquad
  1\le x,\qquad 1\le\log x,\qquad
  \log q\le\log x,\qquad 0<\log q.
\]
\end{helper}
\begin{proof}
The proof uses only standard asymptotics and the dyadic pigeonhole step used
in the paper.  Since \(\log x=o(x)\) and
\(\log x=o(x^{\delta/2})\) as \(x\to\infty\), choose \(X_0\) so large that
for every \(x\ge X_0\),
\[
  1\le x,\qquad 1\le\log x,\qquad
  \log x\le \frac{\delta}{200}x,\qquad
  \log x\le \frac{\delta}{200}x^{\delta/2}.
\]
Fix such an \(x\), and put
\[
  y=\frac{\delta}{100}\frac{x}{\log x}.
\]
The inequality \(\log x\le(\delta/200)x\), together with \(\log x>0\), gives
\(2\le y\).

Choose \(m\ge0\) with
\[
  2^m\le y<2^{m+1}.
\]
This is the usual dyadic choice of a power of two below \(y\).  Since
\(y\ge2\), the case \(m=0\) is impossible, so \(m\ge1\).  Let \(q=2^m\).
The displayed inequality gives \(q\le y\le2q\), which is the asserted
two-sided estimate for \((\delta/100)x/\log x\).  Since \(m\ne0\), the
standard cardinality formula for Galois fields gives
\[
  |\mathrm{GaloisField}(2,m)|=2^m=q.
\]

It remains to record the lower comparison and logarithmic facts.  From
\(\log x\le(\delta/200)x^{\delta/2}\) and \(x>0\), we have
\[
  x^{1-\delta/2}\cdot 2\log x
    \le x^{1-\delta/2}\cdot 2\cdot\frac{\delta}{200}x^{\delta/2}
    = \frac{\delta}{100}x,
\]
because \(x^{1-\delta/2}x^{\delta/2}=x\).  Dividing by \(2\log x>0\) gives
\[
  x^{1-\delta/2}\le \frac{y}{2}.
\]
The dyadic upper bound \(y<2q\) gives \(y/2\le q\), hence
\(x^{1-\delta/2}\le q\).

The inequalities \(1\le x\) and \(1\le\log x\) hold by the choice of \(X_0\).
Also \(\delta/100\le1\), so
\[
  q\le y=\frac{\delta}{100}\frac{x}{\log x}\le \frac{x}{\log x}\le x
\]
because \(\log x\ge1\).  Monotonicity of the logarithm on positive reals gives
\(\log q\le\log x\).  Finally, \(m\ge1\) implies \(q=2^m\ge2>1\), so
\(0<\log q\).
\end{proof}
